\def\level{BTS}  % Classe à droite
\def\course{CIEL 2}
\def\chaptername{Equations différentielles} % Titre au centre
\def\docpart{Résolution ED avec Geogebra} % Partie du cours
\def\docname{C213}        % Identifiant à gauche

\pagestyle{cours_FP}
\makeheaderBTS{} % Affiche le bandeau

On considère l'équation différentielle $(E):\; y'+2y=-5\e^{-2x}$ où $y$ est une fonction de la variable réelle $x$, définie et dérivable sur $\R$ et $y'$ est la fonction dérivée de $y$.

\begin{enumerate}
	\item Résoudre l'équation homogène $(E_0):\; y'+2y=0$.
	
	\begin{blocvisiblex}[height=2]
    {\quad}
L'équation homogène associée à l'équation différentielle $(E)$ est: $y'+2y=0$. On a donc $a=1$ et $b=2$.
\par
La solution de $(E_0)$ est donc de la forme $y_0(x)=k \times e^{-\frac{2}{1}x}=k \times e^{-2x}$ où $k$ est une constante quelconque.
\end{blocvisiblex}

\medskip

	\item A l'aide du logiciel Géogébra, déterminer les nombres réels $a$ et $b$ pour que la fonction $p$ définie sur $\R$ par: $p(x)=a+bx\e^{-2x}$ soit une solution particulière de l'équation différentielle $(E)$ et donner l'expression de $p$.

\begin{blocvisiblex}[height=1]
    {
			On trouve : $a=$\makebox[0.1\linewidth]{\dotfill} \hfill $b=$\makebox[0.1\linewidth]{\dotfill} \hfill $p(x)=$\makebox[0.3\linewidth]{\dotfill}

	}
	On trouve : $a=$\makebox[0.1\linewidth]{\dotfill} \hfill $b=$\makebox[0.1\linewidth]{\dotfill} \hfill $p(x)=$\makebox[0.3\linewidth]{\dotfill}
\end{blocvisiblex}

	
	\begin{methode}[Créer les membres de gauche et de droite de l'ED dans Geogebra]{}
	1) Dans le champs de saisie (en bas) : taper \verb|p(x)=a+b*x*exp(-2x)|
	
	2) mettre la courbe de $p$ en vert
	
	3) Dans le champs de saisie, taper : \verb|g(x)=p'(x)+2p(x)| {\it \footnotesize (la partie gauche de E avec $p$ au lieu de $y$)}
	
	4) mettre la courbe de $g$ en rouge
	
	5) Dans le champs de saisie, taper : \verb|d(x)= -5*exp(-2x)| {\it \footnotesize (la partie droite de E)}
	
	6) mettre $d$ en bleu
	
	7) faire bouger les curseurs (au besoin, changer les paramètres) pour faire correspondre la courbe rouge avec la courbe bleue {\it \footnotesize (donc la partie gauche avec la partie droite de E)}
	\end{methode}
\medskip

	\item En déduire la forme générale des solutions de (E).
	
	
\begin{blocvisiblex}[height=1.5]
    {
\quad
	}
\end{blocvisiblex}

	\medskip

	\item Déterminer, à l'aide de géogebra et d'un curseur, la solution particulière $f$ de (E) telle que $f(0)=1$.
	


\begin{blocvisiblex}[height=1.5]
    {
\quad
	}
\end{blocvisiblex}

	\begin{methode}[Trouver la solution qui vérifie une condition donnée avec Geogebra]{}
	1) Dans le champs de saisie créer une nouvelle fonction ({\bf ne pas} l'appeler $y$)
	
	2) la mettre en orange	
	
	3) bouger le curseur pour répondre à la question 4
	\end{methode}

\end{enumerate}

\medskip

\begin{exercice_cours}[Résolution d'une équation différentielle du second ordre avec Geogebra]{}

On considère l’équation différentielle (E) : $4y' + y = 3e^{0,5x}-1$ où $y$ est une fonction de la variable réelle $x$, définie et dérivable sur l’ensemble $\R$ des nombres réels, et $y'$ est la fonction dérivée de $y$.

\begin{enumerate}
	\item Résoudre l'équation homogène (E$_0$) : $4y'+y=0$

	\item A l’aide du logiciel Géogébra, déterminer les nombres réels $a$ et $b$ pour que la fonction $p$ définie sur $\R$ par : $p(x) = a + b~e^{0,5x}$ soit une solution particulière de l'équation différentielle (E)
	
	\item En déduire la forme générale des solutions de (E).

	\item Déterminer, à l'aide de géogebra et d'un curseur, la solution $f$ de (E) telle que $f(0)=4$.

\end{enumerate}
\end{exercice_cours}